\documentclass{memoir}
\usepackage[utf8]{inputenc}
\usepackage[T1]{fontenc}
\usepackage{geometry}
\usepackage{graphicx}
\usepackage{float}
\usepackage{hyperref}
\usepackage[english,indonesian]{babel}

\title{Optimasi Pemrosesan Template LaTeX Menggunakan Pendekatan Berbasis AI}
\author{Academic Support Agent}
\date{\today}

\begin{document}
\maketitle

\begin{abstract}

\end{abstract}

\chapter{PENDAHULUAN}


\chapter{TINJAUAN PUSTAKA}


% Original content (commented out):
% %%% Modelo de Lista de Exercícios para Universidade
%%% Criado por Diogo Roberto R. Freitas (diogo@poli.br)
%%% Livre para alterações
\documentclass[12pt,onepage,a4paper]{memoir}

%% Language and font encodings
\usepackage[english,portuges]{babel}
\usepackage[T1]{fontenc}
\usepackage[utf8]{inputenc}

%% Sets page size and margins
\usepackage[a4paper,top=2.5cm,bottom=2cm,left=2cm,right=2cm,marginparwidth=1.75cm]{geometry}
\setlength\parindent{0cm} % Tamanho da tabulação dos parágrafos

%%% ATENÇÃO!!! %%%
%%% Preencha estes comandos com suas informações
\newcommand{\logo}{\includegraphics[width=0.4\textwidth]{Logo_poli_upe.png}} % inclua o arquivo com o logo da instituição
\newcommand{\univ}{Universidade de Pernambuco}
\newcommand{\escola}{Escola Politécnica de Pernambuco}
\newcommand{\disc}{Eletrônica Industrial}
\newcommand{\auth}{Prof. Diogo Roberto}
\newcommand{\email}{\url{diogo@poli.br}}
\newcommand{\sitedisc}{\url{canvas.instructure.com/courses/1139439}} % site da disciplina
\newcommand{\cabec}{UPE –-- Escola Politécnica de Pernambuco} % cabeçalho a partir da 2ª página
\newcommand{\tit}{Lista de exercícios 2}
%%%%%%%%%

%% Useful packages
\usepackage{amsmath}
\usepackage{graphicx}
\usepackage[colorinlistoftodos]{todonotes}
\usepackage[colorlinks=true, allcolors=blue]{hyperref}
\linespread{1.25}
\usepackage{graphicx}
\usepackage{nicefrac}
\usepackage[tight]{units}
\usepackage[justification=centering]{caption}
\usepackage{subcaption}
\usepackage{lastpage}
\usepackage{pstricks}
\usepackage{url}% ou hyperref
%\usepackage{breakurl}
\usepackage{multirow}
\usepackage{tabulary}
\usepackage{longtable}
\usepackage{microtype}% improves the spacing between words and letters
\usepackage{booktabs}%  helps you improve the quality of your LaTeX tables
\usepackage{rotfloat}
\usepackage{rotating}

%%% NÃO ALTERAR ESTES COMANDOS 
%%% Configura a primeira página
\makepagestyle{1pagina}
\makeoddhead{1pagina}{
	\logo \\
    \vspace{5pt}
    %\textsf{\univ \\ \escola \\} %%% Caso o logo não tenha o nome da universidade descomente essa linha
    \textsf{Disciplina: \disc \\
	\auth~(\email) \\
    \sitedisc }
    }{}{}
\makeoddfoot{1pagina}{\tiny Documento produzido em \LaTeX}{}{\scriptsize Página \thepage~de \thelastpage}
\makefootrule{1pagina}{\textwidth}{\normalrulethickness}{5pt}
%%% Configura as demais páginas
\makepagestyle{paginacomum}
\makeevenhead{paginacomum}{\textsf{\scriptsize \cabec~-- \auth}}{}{}
\makeevenfoot{paginacomum}{\tiny Documento produzido em \LaTeX}{}{\scriptsize Página \thepage~de \thelastpage}
\makeoddhead{paginacomum}{\textsf{\scriptsize \cabec~-- \auth}}{}{}
\makeoddfoot{paginacomum}{\tiny Documento produzido em \LaTeX}{}{\scriptsize Página \thepage~de \thelastpage}
\makeheadrule{paginacomum}{\textwidth}{\normalrulethickness}
\makefootrule{paginacomum}{\textwidth}{\normalrulethickness}{5pt}
\pagestyle{paginacomum}
%%%


%%% Início do documento
\begin{document}
\thispagestyle{1pagina}
\vspace*{3.5cm} %%% Caso o cabeçalho cubra estes texto aumente o "vspace"
\textbf{\textsf{\large \tit}} %%% não altere aqui

\begin{enumerate} % início das questões, pode escrever aqui
% Questão 1
\item Hi, olá, sãoção. aswb wejw ejw ewje wjeh wjejw ejwh xj ecj ejw jec wjew cje wjch ewjhce wjeh cwjehwcjhe wjh cejweh cwhe wcje cwjh ewjeh wjehc wje wje .
	\begin{enumerate}[a.]
	\item aaaqw mnnmkj bhh.
    \item ahggef aggf.
    \item ção zarção ö
	\end{enumerate}
\newpage
% Questão 2
\item Hi, olá, sãoção. aswb wejw ejw ewje wjeh wjejw ejwh xj ecj ejw jec wjew cjelá, sãoção. aswb wejw ejw ewje wjeh wjejw ejwh xj ecj ejw jec wjew cje wjch  wjch 

% Questão 3
\item Hi, olá, sãoção. aswb wejw ejw ewje wjeh wjejw ejwh xj ecj ejw jec wjew cje wjch lá, sãoção. aswb wejw ejw ewje wjeh wjejw ejwh xj ecj ejw jec wjew cje wjch 11111

\end{enumerate} % fim das questões
\end{document}

\end{document}